\documentclass[11pt]{article}
\usepackage[margin=1in]{geometry}
\usepackage{amsmath,amssymb,amsthm,mathtools}
\usepackage[T1]{fontenc}
\usepackage{lmodern}
\usepackage{microtype}
\usepackage[colorlinks=true,linkcolor=blue,citecolor=blue,urlcolor=blue]{hyperref}
\hypersetup{pdftitle={Exact Maximal Dynamics of the Normalized Inverse-Weight Lin--Lu--Yau Flow on a Weighted Triangle},pdfauthor={Xiaojian Zeng}}

\newtheorem{theorem}{Theorem}
\newtheorem{proposition}[theorem]{Proposition}
\newtheorem{lemma}[theorem]{Lemma}
\theoremstyle{remark}
\newtheorem{remark}[theorem]{Remark}

\newcommand{\G}{\mathcal G}
\newcommand{\D}{\Delta_3}
\newcommand{\R}{\mathbb R}
\DeclareMathOperator{\Lip}{Lip}

\title{Exact Maximal Dynamics of the Normalized Inverse-Weight\\
Lin--Lu--Yau Flow on a Weighted Triangle}
\author{Xiaojian Zeng\\
\small PBC School of Finance, Tsinghua University, 100190\\
\small \texttt{xiaojianzeng@foxmail.com}}
\date{}

\begin{document}
\maketitle

\begin{abstract}
We completely classify the normalized inverse-weight Lin--Lu--Yau flow on a
labelled weighted triangle when the curvature quotient is divided literally
by the edge length.  For positive lengths $(a,b,c)$ summing to one, all three
curvatures are finite exactly when $\max\{a,b,c\}\leq 1/2$.  On this region
we give exact primal and dual transport certificates, including every
shortest-path and optimal-transport equality wall, and obtain a Lipschitz
piecewise-rational vector field.  Every admissible initial value has a unique
global classical trajectory and no trajectory reaches a new transport wall,
the non-finiteness locus, or a zero edge in finite time.  After sorting
$a\geq b\geq c$, the equal point is fixed; data with $a>b=c$ converge to
$(1/2,1/4,1/4)$; and all data with $b>c$ converge to
$(1/2,1/2,0)$.  These alternatives give all equilibria, boundary
degenerations, and omega-limit sets.  The last limit belongs only to the
closed simplex; no curvature is assigned at its zero edge.
\end{abstract}

\section{Setting and relation to earlier flows}

Ollivier's coarse Ricci curvature compares probability measures in
Wasserstein distance, and Lin, Lu, and Yau introduced the graph idleness
derivative that now bears their names
\cite{Ollivier2009,LinLuYau2011}.  Bai, Lin, Lu, Wang, and Yau developed a
continuous normalized flow for weighted graphs
\cite{BaiLinLuWangYau2024}.  Their curvature is normalized by the induced
endpoint distance $d_w(x,y)$, and their flow procedure deletes an edge once
its length ceases to realize that distance.  Bai and Hua subsequently studied
the inverse-weight specialization on finite trees, where every edge is
automatically geodesic, and recorded its normalized flow and omega-limit
notation \cite{BaiHua2026}.  They also noted that the behavior of this
inverse-weight flow on graphs with cycles remained untreated.

Here we solve the literal no-surgery triangle problem.  The distinction is
essential: the denominator below is the edge length even when that edge is
not a shortest path.  The result is confined to $C_3$ and does not import the
tree curvature formula, tree Gauss--Bonnet identity, or tree convergence
theorems.

Let the edges $e_{12},e_{23},e_{31}$ have respective lengths
\[
 a,b,c>0,\qquad a+b+c=1,
\]
and let $d_w$ be the induced shortest-path metric.  At a vertex $x$, set
\[
 D_x=\sum_{e\ni x}\frac1{w_e},
\]
and, for $0\leq\alpha<1$, define
\[
 m_x^\alpha(x)=\alpha,\qquad
 m_x^\alpha(z)=(1-\alpha)\frac{w_{xz}^{-1}}{D_x}\quad(z\sim x).
\]
For $e=xy$ of length $w_e$, the curvature under study is
\begin{equation}\label{eq:literal-curvature}
 \kappa_e(w)=\lim_{\alpha\uparrow1}
 \frac{1-W_1(m_x^\alpha,m_y^\alpha)/w_e}{1-\alpha},
\end{equation}
whenever this is a finite real number.  When all three curvatures are finite,
put
\begin{equation}\label{eq:flow}
 K=a\kappa_a+b\kappa_b+c\kappa_c,
 \qquad \dot x=-x(\kappa_x-K),\quad x\in\{a,b,c\}.
\end{equation}
We write
\[
 \D=\{(a,b,c)>0:a+b+c=1\},\qquad
 \G=\{(a,b,c)\in\D:\max(a,b,c)\leq\tfrac12\}.
\]
For a global trajectory, its omega-limit set means
\[
 \omega(w)=\{w^*\in\overline\D:\text{there are }t_j\to\infty
 \text{ with }w(t_j)\to w^*\}.
\]

\section{Exact transport cells}

We derive the curvature of $e_{12}$; cyclic relabelling gives the other two.
Put $\varepsilon=1-\alpha$, $\mu=m_1^{1-\varepsilon}$, and
$\nu=m_2^{1-\varepsilon}$.  In vertex order $(v_1,v_2,v_3)$,
\[
 \mu=\left(1-\varepsilon,\frac{\varepsilon c}{a+c},
                    \frac{\varepsilon a}{a+c}\right),\qquad
 \nu=\left(\frac{\varepsilon b}{a+b},1-\varepsilon,
                    \frac{\varepsilon a}{a+b}\right).
\]
Thus $\Delta=\mu-\nu$ has coordinates
\begin{align}
 \Delta_1&=1-\varepsilon-\frac{\varepsilon b}{a+b},\nonumber\\
 \Delta_2&=-1+\varepsilon+\frac{\varepsilon c}{a+c},\label{eq:imbalance}\\
 \Delta_3&=\frac{\varepsilon a(b-c)}{(a+b)(a+c)},
 \qquad \Delta_1+\Delta_2+\Delta_3=0.\nonumber
\end{align}
For
\[
0<\varepsilon<\min\left\{\frac{a+b}{a+2b},
                           \frac{a+c}{a+2c}\right\},
\]
we have $\Delta_1>0$, $\Delta_2<0$, and
$\operatorname{sgn}\Delta_3=\operatorname{sgn}(b-c)$.

Write
\[
 A=d_w(v_1,v_2),\quad B=d_w(v_2,v_3),\quad C=d_w(v_3,v_1).
\]
If $a\leq1/2$, then $A=a$.  We use the finite transport primal and
Lipschitz-potential dual recalled in \cite[Section 2]{BaiLinLuWangYau2024}.
We now certify the optimum explicitly.  If $b\geq c$, the
nonzero entries
\[
 \pi_{11}=\nu_1,\quad \pi_{22}=\mu_2,\quad \pi_{33}=\nu_3,
 \quad \pi_{12}=\Delta_1,\quad \pi_{32}=\Delta_3
\]
form a coupling of cost $a\Delta_1+B\Delta_3$.  The Kantorovich potential
\[
 (\phi_1,\phi_2,\phi_3)=(a,0,B)
\]
has the same dual value and is $1$-Lipschitz: its only non-immediate
constraint is $|a-B|\leq C$, which is a metric triangle inequality.  Hence
primal and dual values coincide.  If $c\geq b$, instead use
\[
 \pi_{11}=\nu_1,\quad \pi_{22}=\mu_2,\quad \pi_{33}=\mu_3,
 \quad \pi_{12}=-\Delta_2,\quad \pi_{13}=-\Delta_3.
\]
Its cost is $a(-\Delta_2)+C(-\Delta_3)$, attained in the dual by
$(a,0,a-C)$; feasibility follows from $|a-C|\leq B$.  At $b=c$ the residual
third-vertex shipment vanishes, so the certificates include the transport
equality wall.

On $\G$ we have $B=b$ and $C=c$.  Substitution in
\eqref{eq:literal-curvature} gives
\begin{equation}\label{eq:two-cells}
 \kappa_a=
 \begin{cases}
  \displaystyle 1+\frac{b(a+2c-b)}{(a+b)(a+c)},&b\geq c,\\[6pt]
  \displaystyle 1+\frac{c(a+2b-c)}{(a+b)(a+c)},&c\geq b.
 \end{cases}
\end{equation}
The two expressions differ by
\[
 \frac{(b-c)(a-b-c)}{(a+b)(a+c)},
\]
and therefore agree when $b=c$.  More symmetrically, if $x$ is any edge and
$y,z$ are the other two, let $M=\max(y,z)$ and $m=\min(y,z)$.  Cycling
\eqref{eq:two-cells} yields the complete simultaneous transport-cell formula
\begin{equation}\label{eq:curvature-vector}
 \displaystyle
 \kappa_x=1+\frac{M(x+2m-M)}{(x+y)(x+z)}\qquad(w\in\G).
\end{equation}
Thus the transport partition consists of the six weak ordering chambers and
their equality strata, refined by the three geodesic walls.

The geodesic walls and the non-geodesic region behave differently.  If, for
example, $a=1/2$, then $b+c=1/2$ and either branch gives
\begin{equation}\label{eq:tied-edge}
 \kappa_a=1+\frac{3bc}{(a+b)(a+c)},
\end{equation}
so the shortest-path tie causes no divergence.  If $a>1/2$, however,
$m_1^{1-\varepsilon}\to\delta_{v_1}$ and
$m_2^{1-\varepsilon}\to\delta_{v_2}$, while continuity of $W_1$ on this
finite metric space gives
\[
 W_1(m_1^{1-\varepsilon},m_2^{1-\varepsilon})\longrightarrow
 d_w(v_1,v_2)=b+c=1-a.
\]
The numerator of \eqref{eq:literal-curvature} consequently tends to
$(2a-1)/a>0$, so the quotient tends to $+\infty$.  Cyclically, all three
curvatures are finite exactly on $\G$.

\section{Well-posedness and invariant strata}

Equation \eqref{eq:curvature-vector}, together with \eqref{eq:flow}, is an
exact piecewise-rational vector field on every cell.  For later analysis, on
the representative chamber $a\geq b\geq c$ it reads
\begin{align}
 \kappa_a&=1+\frac{b(a+2c-b)}{(a+b)(a+c)},\nonumber\\
 \kappa_b&=1+\frac{a(b+2c-a)}{(a+b)(b+c)},\label{eq:ordered-kappa}\\
 \kappa_c&=1+\frac{a(c+2b-a)}{(a+c)(b+c)},\nonumber
\end{align}
with $K=a\kappa_a+b\kappa_b+c\kappa_c$ and
$F_x=-x(\kappa_x-K)$.  This also makes
\[
 F_a+F_b+F_c=-K+K(a+b+c)=0,
\]
so the simplex sum is preserved.

Let
\[
 \overline\G=\{(a,b,c)\geq0:a+b+c=1,\ \max(a,b,c)\leq\tfrac12\}.
\]
The right side of \eqref{eq:curvature-vector} has a locally Lipschitz
extension to a neighborhood of $\overline\G$: min and max are Lipschitz, and
each pair sum is at least $1/2$ on $\overline\G$.  This extension is used
only for the ODE argument and does not define the original inverse-weight
measure at a zero edge.
The local existence and uniqueness theorem for a locally Lipschitz field
and its compact-set continuation criterion apply to this extension
\cite[Theorem 2.2 and Corollary 2.15]{Teschl2012}.

Every ordering wall is invariant.  Indeed, the vector field is equivariant
under edge permutations, so equal coordinates have equal derivatives.
For a geodesic wall, take $a=1/2$, $b+c=1/2$, and, without loss, $b\geq c$.
The exact formulas are
\[
 \kappa_a=1+\frac{3bc}{(a+b)(a+c)},\qquad
 \kappa_b=1+\frac{c}{a+b},\qquad
 \kappa_c=1+\frac{b}{a+c},
\]
and direct simplification gives
\begin{equation}\label{eq:wall-identity}
 b(\kappa_b-\kappa_a)+c(\kappa_c-\kappa_a)=0.
\end{equation}
Since $a=b+c=1/2$, \eqref{eq:wall-identity} implies $K=\kappa_a$ and hence
$F_a=0$.  All three geodesic walls are therefore invariant.

Local uniqueness prevents a trajectory starting off one of these invariant
walls from first meeting it at a finite time.  Moreover
\[
 \dot x=x(K-\kappa_x)
\]
has a bounded multiplier on $\overline\G$, so a positive coordinate cannot
vanish in finite time.  A solution beginning in $\G$ consequently remains
positive and in $\G$ for every finite time.  Compactness of $\overline\G$
and the local extension give continuation for all $t\geq0$.

We have proved the following well-posedness statement.

\begin{proposition}\label{prop:wellposed}
The full set of initial data for which all three literal curvatures are finite
and \eqref{eq:flow} defines a classical forward flow is $\G$.  Every datum in
$\G$ has a unique global classical trajectory.  The weak ordering chambers,
their equality intersections, the geodesic walls, and their intersections
are invariant strata.  A strict trajectory reaches neither a new transport
wall nor the non-finiteness locus in finite time, and no positive coordinate
vanishes in finite time.
\end{proposition}

\section{Complete maximal dynamics}

On $a\geq b\geq c$, subtraction of \eqref{eq:ordered-kappa} gives the two
exact log-ratio identities
\begin{align}
 \frac{d}{dt}\log\frac ab
 &=\frac{(a-b)(2ab+3c-1)}{(a+b)(a+c)(b+c)},\label{eq:logab}\\
 \frac{d}{dt}\log\frac bc
 &=\frac{2a(b-c)}{(a+b)(a+c)}.\label{eq:logbc}
\end{align}
The largest-coordinate equation also factors as
\begin{equation}\label{eq:Fa}
 F_a=\frac{a(1-2a)(a^2+2ab+3b^2-2b)}{(a+b)(a+c)}.
\end{equation}
All other chambers follow by permuting labels.

At an equilibrium, positivity forces all three curvatures to equal $K$.
Equation \eqref{eq:logbc} first gives $b=c=s$.  Then $a=1-2s$ and
\begin{equation}\label{eq:eqfactor}
 \kappa_b-\kappa_a=
 \frac{(1-3s)(4s-1)}{2s(1-s)},\qquad \frac14\leq s\leq\frac13.
\end{equation}
Hence the only equilibria are
\begin{equation}\label{eq:equilibria}
 (\tfrac13,\tfrac13,\tfrac13)
 \quad\text{and the three labelled permutations of}\quad
 (\tfrac12,\tfrac14,\tfrac14).
\end{equation}

We now classify every trajectory.  Permute its labels once so that
$a(0)\geq b(0)\geq c(0)$; Proposition~\ref{prop:wellposed} preserves this
ordering for every finite time.

If $a=b=c$, the trajectory is the constant equal equilibrium.  Suppose next
that $a>b=c=s$.  This median is invariant, and the simplex relation reduces
the flow exactly to
\begin{equation}\label{eq:median}
 \dot s=-\frac{(1-2s)(1-3s)(4s-1)}{2(1-s)}.
\end{equation}
For $1/4<s<1/3$ the right side is negative.  Thus $s$ decreases, and its
only possible limiting zero of the right side below the initial value is
$1/4$; otherwise $\dot s$ would remain bounded away from zero.  Therefore
\[
 (a,b,c)\longrightarrow(\tfrac12,\tfrac14,\tfrac14).
\]
The endpoint $s=1/4$ is already the equilibrium in
\eqref{eq:equilibria}.

It remains to consider $b>c$, allowing $a=b$.  Let $R=\log(b/c)$.  By
\eqref{eq:logbc}, $\dot R>0$.  In fact $R\to\infty$.  If it were bounded,
then with $r_0=b(0)/c(0)>1$, the inequalities $b/c\geq r_0$ and
$b+c\geq1/2$ would imply
\[
 b-c=\frac{b/c-1}{b/c+1}(b+c)
 \geq\frac{r_0-1}{2(r_0+1)}>0.
\]
Also $a\geq1/3$ and $(a+b)(a+c)\leq1$, so \eqref{eq:logbc} would bound
$\dot R$ below by a positive constant, a contradiction.  Thus $b/c\to
\infty$.  Since $b\leq1/2$, it follows that $c\to0$.  Finally
$a+b\to1$ while $a,b\leq1/2$, and hence
\[
 (a,b,c)\longrightarrow(\tfrac12,\tfrac12,0).
\]

Combining these cases proves the main result.

\begin{theorem}\label{thm:classification}
For the literal normalized inverse-weight Lin--Lu--Yau flow on the weighted
triangle, the maximal classical dynamics are defined exactly for initial data
in $\G$, and every such trajectory is global.  After sorting the initial
coordinates as $a\geq b\geq c$, precisely one of the following occurs:
\begin{enumerate}
 \item $a=b=c$, and the trajectory is constant;
 \item $a>b=c$, and the trajectory converges to
 $(1/2,1/4,1/4)$;
 \item $b>c$ (possibly $a=b$), and the trajectory converges to
 $(1/2,1/2,0)$.
\end{enumerate}
Together with label permutations, \eqref{eq:equilibria} is the complete
equilibrium list.  Every omega-limit set is the singleton consisting of the
corresponding displayed limit.  All degeneration takes place only as
$t\to\infty$; the point $(1/2,1/2,0)$ belongs to the closed simplex and is
used only as a limit.
\end{theorem}

\section{Exact checking and assistance disclosure}

The transport calculation and the chamber identities were additionally
replayed in exact rational arithmetic with two SymPy scripts
\cite{10.7717/peerj-cs.103}.  They verify
the imbalance sum, both primal--dual objective identities, all equality-wall
specializations, simplex preservation, the log-ratio factorizations,
\eqref{eq:wall-identity}, and \eqref{eq:median}.  The proof above supplies the
optimality, sign, invariance, and limiting arguments; successful script
execution is not used as a substitute for those arguments.

For comparison, Zeng's split-graph study separates a finite mathematical
reduction from explicit cover and packing certificates, which an independent
verifier checks \cite{ZengTuza2026}.  Here the paired certificates are a
transport coupling and a Lipschitz potential with equal objective values;
the rational identities are then checked separately.  The common point is
the separation of certificate construction from verification, rather than a
shared curvature theorem.

OpenAI Codex assisted with route exploration, exact-algebra scripts,
adversarial checking, literature screening, proof organization, language
editing, and LaTeX preparation.  The calculations used SymPy with symbolic
rational expressions; no floating-point experiment is evidence for the
theorem.  This is not a proof-assistant formalization.  The author is
responsible for the mathematical statements, executable artifacts,
citations, and final manuscript.

\bibliographystyle{plain}
\bibliography{references}

\end{document}
